\chapter{Research Methodology}\label{ch:method}

This chapter explains how the study was carried out and why each methodological choice was made. It moves
from the research philosophy and design, through the definition of the analytical problems and the data, to
the preprocessing, feature engineering, evaluation and statistical procedures. It closes with the tools used,
the steps taken to make the results valid and reliable, and the ethical considerations.

\section{Research approach}
The study follows a \emph{design-science} and \emph{quantitative} approach. Its main output is an artefact, the
analytics framework, which is built to solve a stated business problem and is then evaluated rigorously against
explicit criteria. This matches the guidelines of \citet{hevner2004} for design-science research in information
systems: the artefact must address a relevant problem, be evaluated rigorously, and be communicated to both
technical and managerial audiences. The empirical analysis is quantitative and uses a large secondary dataset of sales
transactions. The overall lifecycle follows CRISP-DM \citep{shearer2000} as refined into ten phases in the
project proposal (\cref{fig:lifecycle}).

\begin{figure}[htbp]\centering
\begin{tikzpicture}[node distance=4mm and 4mm, every node/.style={font=\small},
  box/.style={draw=brand, rounded corners, fill=brand!8, minimum width=3.3cm, minimum height=9mm, align=center},
  arr/.style={-{Stealth[length=2mm]}, thick, brand}]
\node[box] (p1) {1. Problem\\definition};
\node[box, right=of p1] (p2) {2. Data\\collection};
\node[box, right=of p2] (p3) {3. Data\\preparation};
\node[box, right=of p3] (p4) {4. Exploratory\\analysis};
\node[box, below=9mm of p4] (p5) {5. Feature\\engineering};
\node[box, left=of p5] (p6) {6. Model\\development};
\node[box, left=of p6] (p7) {7. Model\\evaluation};
\node[box, left=of p7] (p8) {8. Dashboard\\development};
\node[box, below=9mm of p8] (p9) {9. System\\evaluation};
\node[box, right=of p9] (p10) {10. Documentation};
\draw[arr] (p1)--(p2); \draw[arr] (p2)--(p3); \draw[arr] (p3)--(p4); \draw[arr] (p4)--(p5);
\draw[arr] (p5)--(p6); \draw[arr] (p6)--(p7); \draw[arr] (p7)--(p8); \draw[arr] (p8)--(p9); \draw[arr] (p9)--(p10);
\draw[arr, dashed, accent] (p7.north) to[bend left=20] node[above, font=\scriptsize]{iterate} (p3.south);
\end{tikzpicture}
\caption{Project lifecycle, adapted from CRISP-DM}\label{fig:lifecycle}
\end{figure}

\section{Research philosophy and design}
\citet{saunders2019} describe research design as a set of layered choices, often drawn as a ``research
onion'', that run from the underlying philosophy to the concrete techniques of data collection and analysis.
\Cref{tab:onion} records the choice made at each layer and the reason for it.

\begin{table}[htbp]\centering\small
\caption{Research design choices, following the layers of \citet{saunders2019}}\label{tab:onion}
\begin{tabularx}{\linewidth}{@{}p{3.1cm}p{3.4cm}X@{}}
\toprule
\textbf{Layer} & \textbf{Choice} & \textbf{Justification} \\ \midrule
Philosophy & Positivism with a pragmatic orientation & Sales outcomes are measurable and the hypotheses can be tested objectively; the pragmatic orientation reflects the aim of building a useful artefact \\
Approach to theory & Deductive & Hypotheses H1--H7 are derived from the literature and then tested on data \\
Methodological choice & Mono-method quantitative & All evidence is numerical: transactions, metrics and test statistics \citep{creswell2018} \\
Strategy & Design-science research with a single in-depth case & An artefact is designed, built and evaluated on one company's data \citep{hevner2004, peffers2007} \\
Time horizon & Longitudinal & 81 months of transactions allow trends, seasonality and out-of-time testing \\
Data collection & Secondary, archival & A transaction export; no human participants are involved \\
Analysis & Descriptive, inferential and predictive & Descriptive statistics, non-parametric tests, econometric and machine-learning models \\ \bottomrule
\end{tabularx}
\end{table}

Within design science, the study follows the six activities of \citet{peffers2007}: problem identification
(\cref{ch:intro}), definition of the objectives of a solution (\cref{ch:intro}), design and development
(\cref{ch:system}), demonstration on the Dan Cake data (\cref{ch:results}), evaluation (\cref{ch:results}) and
communication (this report, the website and the dashboard). The quantitative analysis inside the artefact
follows the CRISP-DM lifecycle shown in \cref{fig:lifecycle}, so the two frameworks operate at different levels:
design science governs the research as a whole, CRISP-DM the analytical work within it.

\section{Problem definition}
Four analytical problems were derived from the business questions:
\begin{enumerate}
  \item \textbf{Demand forecasting (regression).} Predict the number of units of each product sold in the
        next month. This supports production and inventory planning.
  \item \textbf{Large-order prediction (classification).} At the moment an order is entered, predict whether
        it will be large (\(\geq 10\) units). This supports allocation of stock and delivery capacity.
  \item \textbf{Segmentation (clustering).} Group products and market cells (region $\times$ channel $\times$
        customer type) with similar behaviour, to support portfolio and sales-force strategy.
  \item \textbf{Anomaly detection.} Detect unusual trading days and suspicious transactions. This supports
        monitoring and data-quality control.
\end{enumerate}

\section{Data}
The dataset comprises \NRows transactions with \NCols attributes, recorded between \DateMin and \DateMax. Each
record holds an identifier and date; the product (identifier, name, category, pack size, launch year); price
information (printed maximum retail price, list price, market price, discount percentage and unit selling
price); the quantity; the gross, discount and net sales values in Bangladeshi Taka (BDT); and the region,
channel and customer type. \Cref{app:dictionary} gives the full data dictionary. The
data contain no personal identifiers.

\subsection{Dataset profile and validation}
\Cref{tab:dqsummary} summarises the dataset and the results of the validation layer. Each transaction belongs to
one of \NProducts products in 13 categories. These are sold across \NRegions regions through \NChannels
channels (General Trade, Modern Trade, Distributor, E-commerce) to \NCustTypes customer types (Retailer, Super
Shop, Distributor, Corporate, Online Customer). Five products were launched after 2020: Lachcha Semai, two
breads, a croissant and a cappuccino muffin gift pack.

\begin{table}[htbp]\centering\small
\caption{Dataset profile and data-quality validation results}\label{tab:dqsummary}
\begin{tabular}{@{}lr@{}}
\toprule
\textbf{Check} & \textbf{Result} \\ \midrule
Rows $\times$ columns & \NRows $\times$ \NCols \\
Period covered & \DateMin -- \DateMax \\
Products / regions / channels / customer types & \NProducts / \NRegions / \NChannels / \NCustTypes \\
Missing cells (incl.\ hidden placeholders such as ``NA'', ``-'') & 0 \\
Exact duplicate rows / duplicate transaction IDs & 0 / 0 \\
Domain-rule violations (quantity, price, discount, value) & 0 \\
Gross $=$ quantity $\times$ market price (max.\ abs.\ error) & $<10^{-11}$ BDT \\
Net $=$ gross $-$ discount (max.\ abs.\ error) & $<10^{-11}$ BDT \\
Year/month columns consistent with date & 100\% \\
Product ID $\rightarrow$ one name, category, pack size & 100\% \\
Statistical quantity outliers (flagged, retained) & \NOutlierFlags (\OutlierPct\%) \\
Warehouse reconciliation (SQL total $-$ pandas total) & 0.00 BDT \\ \bottomrule
\end{tabular}
\end{table}

\section{Preprocessing pipeline}\label{sec:preprocessing}
The pipeline is written defensively. It must work correctly on dirty data, not only on the delivered file,
which proved to be exceptionally clean. Its behaviour on dirty data was verified in two ways. Six hands-on
practice scripts deliberately damaged copies of the data and measured how much was repaired. In addition, the
automated tests in \texttt{tests/} inject missing values, duplicates and impossible values and assert that the
pipeline repairs them (\cref{sec:testing}). The steps are applied in a fixed order, because some steps depend
on others.

\subsection{Data type conversion and validation}
Dates are parsed with an explicit format (\texttt{dd-mm-yy}) and numeric columns are coerced. Values that fail
to parse become missing, so they are never silently misread. Text fields are trimmed of spaces. The validation
step then produces a machine-readable quality report (\texttt{outputs/metrics/data\_quality.json}). The report
checks for missing values, hidden placeholder strings, duplicate rows and keys, and domain-rule violations. It
also checks the accounting identities between columns and the consistency of calendar fields and product
attributes.

\subsection{Missing-value treatment}
A decision hierarchy was implemented, following \citet{garcia2015}:
\begin{enumerate}
  \item \emph{derive} the value exactly from other columns where an identity exists (for example,
        $\text{quantity} = \text{gross}/\text{market price}$, or year and month from the date);
  \item \emph{drop} the row if an essential identifier (transaction ID, date, product) is missing, because a
        guess would amount to fabrication;
  \item \emph{impute} numeric gaps with the median of the same product, which is robust to skew;
  \item \emph{impute} product attributes with the product's mode, and label missing market attributes
        explicitly as ``Unknown'' so that they remain visible.
\end{enumerate}
In practice tests with 3\% of cells removed at random, deriving values from accounting identities restored
quantities, prices and sales values exactly. Simple mean imputation, by contrast, introduced errors of several
units per order.

\subsection{Duplicate removal}
Exact duplicates are removed first. Duplicate keys, meaning the same transaction ID with differing values, are
resolved by keeping the most complete version. IDs are compared after upper-casing and trimming, so that hidden
duplicates such as `` dc2020000001'' are caught. Near-duplicates (same date, product, region, channel and
quantity but different IDs) are only \emph{flagged}. In sales data they are usually genuine repeat orders.

\subsection{Outlier detection and treatment}
Order quantities are strongly right-skewed (skewness \QtySkew; median \QtyMedian units). Each product also has
its own typical order size. A global $z$-score therefore flags too many small-product orders and too few
large-product ones. Four methods were compared in the practice phase: the IQR rule, the $z$-score, the MAD-based
modified $z$-score, and a per-product IQR rule on the log scale. The last of these was the most balanced, and
it is the method used in the pipeline:
\[
  \text{flag}_i = \bigl[\, \log(1+q_i) \notin [\,Q_1^{(p)} - 1.5\,\text{IQR}^{(p)},\; Q_3^{(p)} + 1.5\,\text{IQR}^{(p)}\,] \bigr],
\]
where the quartiles are computed separately for each product $p$. Impossible values (for example, a quantity
below one or a discount above 50\%) are treated as errors: they are blanked and re-derived. Unusual but possible
values are flagged and kept. The real data contained \NOutlierFlags such flags (\OutlierPct\% of orders), all
plausible bulk orders, and no impossible values.

\subsection{Normalisation and encoding}
Scaling is not applied globally. It is applied \emph{inside} each model's pipeline, fitted on the training rows
only, to avoid leakage. Count-like features are log-transformed, $\log(1+x)$, to stabilise variance, and they
are then standardised for scale-sensitive models (linear regression, KNN and SVR). Nominal categories (region,
channel, customer type, category) are one-hot encoded. Month is encoded cyclically as
$(\sin 2\pi m/12, \cos 2\pi m/12)$, so that December and January are neighbours. High-cardinality interactions
(channel $\times$ region, product) are target-encoded with additive smoothing,
\[
  \tilde\mu_k = \frac{n_k \bar y_k + m\,\bar y}{n_k + m}, \qquad m=50,
\]
using training data only. Rare levels are thereby shrunk towards the overall mean.

\subsection{Feature engineering for demand forecasting}
Transactions were aggregated into a \emph{panel} of \PanelSeries products $\times$ months. Each feature for
month $t$ uses only information available at the end of month $t-1$, or information that is planned in advance
(\cref{tab:features}). The same feature function is used for training, testing and future forecasting, so the
rule is enforced in one place. One instructive check compared a correct three-month rolling mean (shifted
before rolling) with a ``leaky'' version that included the current month. The leaky version correlated more
strongly with the target, which is exactly why leakage is tempting and dangerous \citep{kaufman2012}.

\begin{table}[htbp]\centering\small
\caption{Features of the demand-forecasting model (\PanelFeatures in total)}\label{tab:features}
\begin{tabularx}{\linewidth}{@{}lX@{}}
\toprule
\textbf{Group} & \textbf{Features} \\ \midrule
History (lags) & units in months $t-1, t-2, t-3, t-6, t-12$ \\
Rolling statistics & mean of the previous 3, 6 and 12 months; standard deviation of the previous 3 months \\
Momentum & last month's year-on-year growth; 12-month trend ratio \\
Price (planned) & list price of month $t$; year-on-year list-price change \\
Product & category (one-hot); pack weight in grams; gift-pack indicator; months since launch \\
Calendar & month number; cyclical month (sine, cosine); December and February indicators; days in month; time index \\ \bottomrule
\end{tabularx}
\end{table}

\section{Evaluation design}
\subsection{Temporal hold-out}
Because the problems involve time, every supervised model is trained on the years \textbf{2020--2024} and
tested on \textbf{January 2025 -- September 2026}. This unseen period of \TestMonths months plays the role of
``the future''. Random train/test splits were deliberately not used, because they would let the models learn
from observations that come after the ones they predict \citep{bergmeir2012}.

\subsection{Hyper-parameter tuning}
Hyper-parameters were tuned only on the training years, using an \emph{expanding-window} validation. For
validation year $y \in \{2022, 2023, 2024\}$, the model is fitted on all months before $y$ and scored on $y$.
The configuration with the best mean validation score is then refitted on 2020--2024 and evaluated once on the
test period. For classification, the probability threshold was chosen in the same way: the model was fitted on
2020--2023 and the threshold that maximised F1 on 2024 was kept.

\subsection{Metrics}
For regression, the following measures were used, where $y_i$ is the actual and $\hat y_i$ the predicted value
for $n$ product-months:
\begin{align*}
\text{MAE} &= \tfrac{1}{n}\textstyle\sum_i |y_i-\hat y_i|, &
\text{RMSE} &= \sqrt{\tfrac{1}{n}\textstyle\sum_i (y_i-\hat y_i)^2}, &
R^2 &= 1-\frac{\sum_i (y_i-\hat y_i)^2}{\sum_i (y_i-\bar y)^2},\\
\text{MAPE} &= \tfrac{100}{n}\textstyle\sum_i \frac{|y_i-\hat y_i|}{\max(y_i,1)}, &
\text{WAPE} &= 100\,\frac{\sum_i |y_i-\hat y_i|}{\sum_i y_i}. &&
\end{align*}
WAPE is preferred for business communication because it weights each product by its volume. For
classification, accuracy, precision, recall, F1, ROC-AUC, PR-AUC and the lift in the top decile were reported.
For clustering, the silhouette coefficient, the Davies--Bouldin index and the Calinski--Harabasz index were
reported. Anomaly detectors were evaluated by precision and recall on a copy of the data containing planted
errors of known type.

\subsection{Benchmarks}
Following \citet{makridakis2018}, every forecasting model is compared with two naive rules: \emph{naive}
(next month = this month) and \emph{seasonal naive} (next month = the same month last year). At company level,
Holt--Winters and SARIMA serve as statistical benchmarks. The majority-class rule is the benchmark for
classification.

\section{Statistical analysis}
The exploratory phase used descriptive statistics, Spearman rank correlation (robust to the skewed
distributions found), Kruskal--Wallis tests for differences between groups, and chi-square tests of
independence. Because the sample is very large, almost any difference is statistically significant.
\emph{Effect sizes} ($\varepsilon^2$, $\rho^2$ and Cram\'er's $V$) were therefore reported and used to judge
practical importance. A fixed-effects log--log regression estimated the price elasticity of demand
\citep{wooldridge2020}.

Non-parametric tests were chosen deliberately. Order quantities and values are strongly right-skewed (skewness
\QtySkew and \ValueSkew), so the normality assumption of ANOVA and Pearson correlation is not met
\citep{field2018}. The Kruskal--Wallis statistic compares mean ranks across $k$ groups,
\[
  H = \frac{12}{n(n+1)} \sum_{j=1}^{k} \frac{R_j^2}{n_j} - 3(n+1),
  \qquad
  \varepsilon^2 = \frac{H}{(n^2-1)/(n+1)},
\]
where $R_j$ is the rank sum of group $j$ with $n_j$ members, and $\varepsilon^2$ is the share of rank variance
explained by the grouping. Cram\'er's $V=\sqrt{\chi^2/\bigl(n\,(\min(r,c)-1)\bigr)}$ measures the strength of
association in an $r\times c$ contingency table. Effect sizes were interpreted with the conventions of
\citet{cohen1988}. All tests used a significance level of $\alpha = 0.05$. \Cref{tab:hypmap} links each
hypothesis of \cref{sec:hypotheses} to its test.

\begin{table}[htbp]\centering\small
\caption{Hypotheses, statistical tests and decision rules}\label{tab:hypmap}
\begin{tabularx}{\linewidth}{@{}lXl@{}}
\toprule
\textbf{Hypothesis} & \textbf{Test and evidence} & \textbf{Effect size} \\ \midrule
H1 channel $\rightarrow$ order size & Kruskal--Wallis across 4 channels & $\varepsilon^2$ \\
H2 customer type $\rightarrow$ order size & Kruskal--Wallis across 5 customer types & $\varepsilon^2$ \\
H3 discount $\rightarrow$ order size & Spearman rank correlation & $\rho^2$ \\
H4 monthly seasonality & Kruskal--Wallis on detrended daily sales by month & $\varepsilon^2$ \\
H5 price elasticity & Fixed-effects log--log regression; $t$-test on the price coefficient & elasticity \\
H6 ML beats seasonal naive & Out-of-time RMSE and WAPE on 2025--2026 & \% error reduction \\
H7 channel $\leftrightarrow$ large orders & Chi-square test of independence & Cram\'er's $V$ \\ \bottomrule
\end{tabularx}
\par\smallskip\footnotesize $H_0$ is rejected when $p<0.05$. For H6 the decision rests on the size of the
out-of-time error reduction, which is the quantity of business interest.
\end{table}

\section{Two complementary modelling studies}
The predictive work was carried out twice, with different aims. The \emph{production study}
(\cref{ch:models}) builds the best possible models for the business: rich features, randomised hyper-parameter
search over expanding-window folds, statistical benchmarks and forecast reconciliation. The
\emph{algorithm-level study} (\cref{ch:algorithms}) examines each of the nine algorithms named in the proposal
separately, with a simpler shared feature set, a single 2024 validation year and one transparent script per
algorithm. Its aims are to explain how each algorithm learns and fails, and to check the production results
with an independent implementation.

\section{Tools and environment}
The framework was implemented in Python~3 with pandas and NumPy for data handling \citep{mckinney2022},
SciPy and statsmodels for statistics and time-series models, scikit-learn for machine learning
\citep{pedregosa2011}, XGBoost for gradient boosting \citep{chen2016}, Matplotlib and Plotly for
visualisation, SQLite for the dimensional data store, and Streamlit for the web dashboard. The code is
version-controlled with Git and published on GitHub. The dashboard is deployed on Streamlit Community Cloud and
the website on GitHub Pages. The report was typeset in \LaTeX. All random processes use a fixed seed (42), so every result can be reproduced exactly.

\section{Validity and reliability}
Computational studies face the same four classes of threat as experiments \citep{shadish2002, wohlin2012}.
\Cref{tab:validity} lists the main threats in this study and the measures taken against each.

\begin{table}[htbp]\centering\small
\caption{Threats to validity and how they were mitigated}\label{tab:validity}
\begin{tabularx}{\linewidth}{@{}p{2.6cm}XX@{}}
\toprule
\textbf{Type} & \textbf{Threat} & \textbf{Mitigation} \\ \midrule
Internal & Information leakage from the future inflates accuracy & Strict temporal hold-out; history features shifted before rolling; scalers and encoders fitted on training rows only; an automated test checks that no feature uses the current month \\
Internal & Tuning on the test set & Hyper-parameters and thresholds chosen on validation years; the test period is scored once \\
Construct & Metrics that do not reflect business value & WAPE (volume-weighted) for planners, PR-AUC and lift for imbalanced classes, effect sizes beside $p$-values \\
Statistical conclusion & Significance driven only by sample size & Effect sizes and confidence intervals reported; non-parametric tests for skewed data \\
External & Results specific to one company and period & Framework kept configurable and domain-independent; findings framed as a case study; limitations stated in \cref{sec:limitations} \\
Reliability & Results that cannot be reproduced & One command regenerates every number; fixed seed 42; version control; nine automated tests; an independent second implementation (\cref{ch:algorithms}) \\ \bottomrule
\end{tabularx}
\end{table}

A further risk concerns \emph{missing data}. \citet{little2019} distinguish data missing completely at random,
at random, and not at random; only the first two can be imputed without bias. The preprocessing hierarchy
therefore derives values from exact accounting identities first and imputes statistically only as a last
resort, which keeps the treatment valid under all three mechanisms for the derivable columns.

\section{Ethical considerations}
The dataset contains no names, contact details or other personally identifiable information. Customers are
described only by type (for example, \emph{Retailer}) and region. The data were used solely for academic
purposes, in line with the ethical principles stated in the proposal. Model limitations and uncertainties are
reported transparently (\cref{sec:limitations}), so that forecasts are not over-trusted in decision-making.
The public repository contains only aggregated data, and the raw transaction file is kept private.

\section{Chapter summary}
The study is positivist, deductive and quantitative, and it follows a design-science strategy in which the
analytics framework is the artefact. The \NRows transactions were validated, cleaned and transformed with a
defensive, leakage-safe pipeline. Models are evaluated on a strictly later test period against naive and
statistical benchmarks, with metrics chosen for business relevance, and hypotheses are tested with
non-parametric statistics and effect sizes. \Cref{ch:system} now describes how this methodology was turned into
a working system.
